%% ════════════════════════════════════════════════════════════════════════════
%% Jayraj Kiran Ladkat — CV for the AI-for-ITS proposal (team-member section)
%% Style: Fan-Lei-inspired (article class, full-width, clean typography)
%% Compile: pdfLaTeX  ·  Overleaf: all packages preinstalled
%% ════════════════════════════════════════════════════════════════════════════
\documentclass[11pt,a4paper]{article}

%% ── Encoding & fonts ───────────────────────────────────────────────────────
\usepackage[utf8]{inputenc}
\usepackage[T1]{fontenc}
\usepackage{lmodern}                       % Latin Modern — clean serif default

%% ── Page geometry ──────────────────────────────────────────────────────────
\usepackage[top=1.6cm, bottom=2.0cm, left=1.8cm, right=1.8cm]{geometry}

%% ── Typography ─────────────────────────────────────────────────────────────
\usepackage{microtype}                     % subtle kerning improvements
\setlength{\parindent}{0pt}                % no paragraph indent
\setlength{\parskip}{0pt}                  % no inter-paragraph gap

%% ── Section headers (small caps + coloured rule) ───────────────────────────
\usepackage{titlesec}
\setcounter{secnumdepth}{0}                % suppress section numbers

\titleformat{\section}
  {\large\scshape\color{heading}}          % small caps, coloured
  {}                                        % no label
  {0em}                                     % no indent
  {}                                        % no prefix
  [\color{heading}\titlerule]              % horizontal rule beneath
\titlespacing*{\section}{0pt}{1.3ex plus 0.4ex}{0.8ex plus 0.2ex}

%% ── Colours ────────────────────────────────────────────────────────────────
\usepackage{xcolor}
\definecolor{heading}{RGB}{0, 90, 110}     % dark teal (section headers)
\definecolor{linkblue}{RGB}{10, 80, 160}   % subtle navy (URLs)

%% ── Hyperlinks ─────────────────────────────────────────────────────────────
\usepackage{hyperref}
\hypersetup{
  colorlinks  = true,
  urlcolor    = linkblue,
  linkcolor   = black,
  citecolor   = black,
  pdfauthor   = {Jayraj Kiran Ladkat},
  pdftitle    = {Jayraj Kiran Ladkat - CV},
}

%% ── Icons & Graphics ───────────────────────────────────────────────────────
\usepackage{fontawesome5}                  % \faIcon{...} for contact line
\usepackage{graphicx}                      % for figures

%% ── Footer ─────────────────────────────────────────────────────────────────
\usepackage{fancyhdr}
\usepackage{lastpage}
\pagestyle{fancy}
\fancyhf{}
\renewcommand{\headrulewidth}{0pt}
\fancyfoot[L]{\footnotesize\textsc{Last Updated\,:\,Nov.\,5,\,2026}}
\fancyfoot[C]{\footnotesize\textsc{Jayraj Kiran Ladkat\,--\,CV}}
\fancyfoot[R]{\footnotesize\thepage/\pageref{LastPage}}

%% ── Lists ──────────────────────────────────────────────────────────────────
\usepackage{enumitem}
\setlist[itemize]{nosep, leftmargin=1.5em, label=\textbullet, topsep=0.25em}

%% ── Tables ─────────────────────────────────────────────────────────────────
\usepackage{tabularx}


% ════════════════════════════════════════════════════════════════════════════
\begin{document}

% ── Header ──────────────────────────────────────────────────────────────────
\begin{center}
  {\LARGE\bfseries Jayraj Kiran Ladkat}\\[0.5em]
  {\small
    \faIcon{envelope}\;\href{mailto:jayrajkladkat@gmail.com}{jayrajkladkat@gmail.com}%
    \quad
    \faIcon{phone}\;+91\,97170\,61462%
    \quad
    \faIcon{map-marker-alt}\;Gurugram, Haryana, India%
    \quad
    \faIcon{globe}\;\href{https://portfolio-sand-one-17.vercel.app/}{portfolio}%
  }\\[0.25em]
  {\small
    \faIcon{github}\;\href{https://github.com/jayraj-kl}{jayraj-kl}%
    \quad
    \faIcon{linkedin}\;\href{https://www.linkedin.com/in/jayraj-ladkat-25534626b}{jayraj-ladkat}%
  }
\end{center}

\vspace{0.2em}

% ── Summary ─────────────────────────────────────────────────────────────────
\section{Summary}

I build agentic systems and quantize deep learning models. At TIGER-Lab
(University of Waterloo) I am building an iterative harness that improves how
open- and closed-source models write code for rendered video, judged by a
vision-language model (VLM). I co-authored a KL-guided mixed-precision
quantization-aware training method for sub-4-bit inference (IEEE HiPC 2026).
I use VLMs to evaluate generated video and am building up core computer vision
(VLMs, segmentation) for deploying quantized models on edge devices.


% ── Education ───────────────────────────────────────────────────────────────
\section{Education}

\noindent
\begin{tabular*}{\textwidth}{@{}l@{\hspace{1.0em}}l@{\extracolsep{\fill}}r@{}}
  10/2022\,--\,10/2026 & \textbf{B.Tech.\ in Computer Science and Engineering} & MIT World Peace University \\
                       & & Pune, Maharashtra, India \\
\end{tabular*}


% ── Research Interests ──────────────────────────────────────────────────────
\section{Research Interests}

Agentic Systems and Harness Optimisation,\; Model Quantization and Efficient Inference,\;
LLM/VLM Evaluation,\; and AI for Education


% ── Research Experience ─────────────────────────────────────────────────────
\section{Research Experience}

\textbf{Research Intern} \hfill 07/2026\,--\,Present\\
TIGER-Lab, University of Waterloo \hfill \textit{Waterloo, Canada (Remote)}
\begin{itemize}
  \item Advisor: Max Ku.
  \item Building an iterative harness-improvement loop for code-based video
        generation (Manim, HyperFrames) across open- and closed-source models,
        improving on the pi agent harness baseline. Three agents divide the work:
        a \emph{creator} writes the render code, an \emph{explorer} probes the
        harness, and an \emph{evaluator} scores the result.
  \item The explorer runs a genetic algorithm over combinations of prompts and
        user-preference representations, selecting for the widest diversity of
        behaviour changes at each step. It tests how well the harness follows a
        3-tier, 14-dimension preference taxonomy (content, delivery, visual style).
  \item Cut context bloat by 35\% through progressive disclosure, loading
        render-API documentation and tools only when the agent needs them.
  \item Evaluates rendered videos with a persona-conditioned VLM-as-judge.
        Over the pi agent baseline: +11\% render success rate and a +4 judge score.
\end{itemize}

\vspace{0.4em}

\textbf{Quantization Research} \hfill \textit{IEEE HiPC 2026}\\
MIT World Peace University \hfill \textit{Pune, Maharashtra}
\begin{itemize}
  \item Co-authored a co-adaptive mixed-precision QAT method: a budgeted
        controller computes a per-layer marginal KL-divergence signal
        \emph{during} training and re-allocates weight precision (2/4/8-bit,
        per-channel STE fake quantization, int8 activations) while the weights
        are still learning.
  \item At a 3-bit average budget it beats a post-hoc KL allocator by about 3$\times$
        in perplexity on Mamba-130M (314 vs.\ 950), and the advantage transfers
        to GPT-2 (56.4 vs.\ 145.1). The gain holds only below a 4-bit average;
        at 4 bits and above, uniform quantization wins.
\end{itemize}


% ── Publications ────────────────────────────────────────────────────────────
\section{Publications}

\hangindent=1.6em\hangafter=1
P2.\quad Shreya Mandekar, \textbf{Jayraj Kiran Ladkat}, Madhura Phatak,
and Bhavana Tiple,
``Co-Adaptation over Allocation: KL-Guided Mixed-Precision
Quantization-Aware Training,''
\textit{accepted poster},
IEEE International Conference on High Performance Computing,
Data, and Analytics (HiPC), 2026.
\href{https://drive.google.com/file/d/1wYned37BXW-kTc7B7jNQFcwUzh66sxkG/view?usp=sharing}{[\faIcon{external-link-alt}\;Acceptance Proof]}

\vspace{0.45em}

\hangindent=1.6em\hangafter=1
P1.\quad \textbf{Jayraj Kiran Ladkat} and Pranati Waghodekar,
``Real-World Scalability of PostgreSQL: Practical Techniques
Use Case Study,''
in \textit{IEEE International Conference on Blockchain and
Distributed Systems Security} (ICBDS 2025).
\href{https://ieeexplore.ieee.org/document/11377486/}{[\faIcon{external-link-alt}\;IEEE~Xplore~11377486]}


% ── Professional Experience ────────────────────────────────────────────────
\section{Professional Experience}

\textbf{Engineering Intern} \hfill 08/2025\,--\,12/2025\\
Publicis Sapient \hfill \textit{Gurugram, Haryana}
\begin{itemize}
  \item Optimised an AI agent data-generation pipeline: reduced per-run
        token usage, redesigned the help interface, and improved synthetic
        dataset quality.
  \item Extended an AI recruitment screening tool with request
        prefetching, a responsive multi-device frontend, and structured
        assessment logic, cutting evaluation turnaround and manual HR
        review effort.
\end{itemize}

\vspace{0.4em}

\textbf{Software Engineering Intern} \hfill 06/2025\,--\,08/2025\\
Elixir AI \hfill \textit{Pune, Maharashtra}
\begin{itemize}
  \item Built the full-stack geospatial application (Flask, PostgreSQL/PostGIS,
        React, Docker) for the \textit{Climate Smart Public Health Platform},
        a Harvard University--Malagasy Ministry of Public Health collaboration.
\end{itemize}

\vspace{0.4em}

\textbf{Cloud \& Full-Stack Intern} \hfill 01/2025\,--\,02/2025\\
Centre for Development of Advanced Computing (C-DAC) \hfill \textit{Pune, Maharashtra}
\begin{itemize}
  \item Introduced NATS-based asynchronous messaging into the service
        architecture, reducing request wait time under load.
\end{itemize}


% ── Selected Projects ──────────────────────────────────────────────────────
\section{Selected Projects}

\textbf{\href{https://github.com/jayraj-kl/dexterio}{Dexterio}%
\,---\,Autonomous Educational Animation Generation Engine}
\hfill 2026
\begin{itemize}
  \item Engineered an agentic RAG pipeline orchestrating multi-step lesson storyboard planning and modular Manim animation code generation grounded in course materials.
  \item Architected a distributed worker fleet using Apache Kafka, Daytona/Docker remote VM sandboxes for safe code compilation, and PostgreSQL-backed LangGraph checkpointers for stateful thread resumption.
\end{itemize}

\vspace{0.35em}

\textbf{\href{https://github.com/jayraj-kl/stud-ex}{StudEx}%
\,---\,RAG Tutors for Core CS Subjects}
\hfill 2024\,--\,2025
\begin{itemize}
  \item Built an educational platform with subject-specific chatbots (Theory of Computation, ML, Blockchain, Cloud) grounded in course PDFs via OpenAI embeddings and a Pinecone vector store (LangChain orchestration).
\end{itemize}


% ── Technical Skills ────────────────────────────────────────────────────────
\section{Technical Skills}

\begin{tabularx}{\textwidth}{@{}>{\bfseries\small}r@{\hspace{1.2em}}X@{}}
  Agentic systems: & Multi-agent harnesses, LangGraph, genetic/evolutionary prompt
                     search, LLM/VLM-as-judge evaluation, progressive disclosure,
                     agentic RAG (LangChain, Pinecone) \\[0.3em]
  Quantization:    & Quantization-aware training, mixed-precision allocation,
                     STE fake quantization, KL-based sensitivity analysis, PyTorch \\[0.3em]
  Vision:          & VLM-based evaluation of rendered video \\[0.3em]
  Languages:       & Python, TypeScript, JavaScript, SQL, C/C++ \\[0.3em]
  Systems:         & Docker, Apache Kafka, NATS, PostgreSQL/PostGIS, AWS (ECS, S3),
                     Flask, Node.js, React, Next.js, Git \\
\end{tabularx}

\end{document}
